\begin{pastexam}[label={px:nccu-114-dsa-06}]{政大資科所 114 學年度 第 6 題}
（單選；每題 4 分）

(4\%) Given the following algorithm, which is likely to be the result if the input variable N is
-1?
\begin{minted}[fontsize=\small, xleftmargin=4.5cm]{text}
Boolean Magic (int N)
{   if (N==1)
      Return True
    else if (N==0)
      Return False
    else { N= N-2;
      Return Magic(N);
        };
}
\end{minted}
A.~-3 \qquad B.~True \qquad C.~False \qquad D.~Stack Overflow \qquad E.~Syntax Error
\end{pastexam}

\begin{pastsoln}
答案是 D。

追蹤呼叫：
\[
\proc{Magic}(-1) \to \proc{Magic}(-3) \to \proc{Magic}(-5) \to \cdots
\]
$N$ 從 $-1$ 開始每次減 $2$，永遠是負的奇數，不會等於 $1$ 也不會等於 $0$，兩個終止條件都碰不到。
這個遞迴只有在 $N$ 是不小於 $0$ 的整數時才會停（偶數停在 $0$、奇數停在 $1$），負數輸入沒有被
處理。

每一次呼叫都要在呼叫堆疊上推一個活動記錄（參數 $N$、返回位址，資結書第~6.4~節），呼叫一直不
返回，堆疊越長越深，很快就超過系統給的堆疊空間，程式以 stack overflow 終止。 （如果 \texttt{int}
是 32 位元，理論上要減到超過 $2^{31}$ 才會溢位繞回正數，那需要大約 $10^9$ 層呼叫，堆疊早就爆了；
何況有號整數溢位在 C 裡是未定義行為。）

逐項檢查：
\begin{itemize}
\item A.~$-3$：函數回傳的是布林值，$-3$ 只是第二層呼叫的參數，錯。
\item B.~True、C.~False：要碰到 $N = 1$ 或 $N = 0$ 才會回傳，負的奇數碰不到，錯。
\item D.~Stack Overflow：無窮遞迴，對。
\item E.~Syntax Error：題目是虛擬碼風格的演算法，問的是執行結果，語法沒有問題，錯。
\end{itemize}
\end{pastsoln}
